\chapter{3-8译码器——5大变式详解}

\section{变式1：74LS138风格译码器（低有效输出+三使能端）——必考！}

\begin{detailbox}[完整教学]

\textbf{【电路】}
74LS138是考试中最常出现的译码器芯片。特点：
\begin{itemize}
  \item 3位地址输入A2,A1,A0
  \item 8个输出Y0-Y7，\textbf{全部低电平有效}（选中=0，未选中=1）
  \item 3个使能端：G1(高有效), G2A(低有效), G2B(低有效)
  \item 只有三个使能端同时有效（G1=1, G2A=0, G2B=0）时，译码器才工作
\end{itemize}

\textbf{【使能逻辑】}$enable = G1 \cdot \overline{G2A} \cdot \overline{G2B}$

\textbf{【VHDL】}
\begin{lstlisting}
entity decoder_74ls138 is
  port (
    A    : in  std_logic_vector(2 downto 0);
    G1   : in  std_logic;
    G2A_n: in  std_logic;  -- 低有效
    G2B_n: in  std_logic;  -- 低有效
    Y_n  : out std_logic_vector(7 downto 0)  -- 低有效
  );
end decoder_74ls138;

architecture behavioral of decoder_74ls138 is
  signal enable : std_logic;
begin
  enable <= G1 and (not G2A_n) and (not G2B_n);

  process(A, enable)
  begin
    if enable = '0' then
      Y_n <= "11111111";  -- 全高=全无效
    else
      case A is
        when "000" => Y_n <= "11111110";  -- Y0=0
        when "001" => Y_n <= "11111101";  -- Y1=0
        when "010" => Y_n <= "11111011";
        when "011" => Y_n <= "11110111";
        when "100" => Y_n <= "11101111";
        when "101" => Y_n <= "11011111";
        when "110" => Y_n <= "10111111";
        when "111" => Y_n <= "01111111";
        when others => Y_n <= "11111111";
      end case;
    end if;
  end process;
end behavioral;
\end{lstlisting}

\textbf{【为什么用3个使能端】}
多个使能端方便级联。G2A和G2B是低有效，G1是高有效——这种混合设计使得可以用不同信号控制，增加灵活性。

\textbf{【常考陷阱】}
\begin{itemize}
  \item 看清输出是低有效还是高有效——74LS138输出0表示选中
  \item 三个使能端必须\textbf{同时}满足条件译码器才工作
  \item 使能无效时所有输出为1（全高=全无效）
\end{itemize}

\end{detailbox}


\section{变式2：用译码器实现任意逻辑函数（必考！）}

\begin{detailbox}[完整教学]

\textbf{【核心原理】}
译码器的每个输出=一个最小项。\textbf{任何组合逻辑函数}=若干最小项之和（OR）。

\textbf{【通用公式】}$F(A,B,C) = \sum m(\dots)= Y_{i1}+Y_{i2}+Y_{i3}+\dots$

\textbf{【例题】}用3-8译码器+一个OR门实现$F(A,B,C)=\sum m(1,3,5,7)$。

\textbf{【解法】}
\begin{enumerate}
  \item 译码器产生8个最小项：Y0=m0, Y1=m1, ..., Y7=m7
  \item $F = m1+m3+m5+m7 = Y1+Y3+Y5+Y7$
  \item 用1个4输入OR门连接Y1,Y3,Y5,Y7
\end{enumerate}

\textbf{【电路结构】}
\begin{center}
\begin{tikzpicture}
\node[draw,minimum width=2.5cm,minimum height=1.5cm] (dec) at (0,0) {3-8译码器};
\node[american or port] (orgate) at (4,0) {};
\draw[->] (dec.east|-orgate.in 1) -- (orgate.in 1);
\draw[->] (dec.east|-orgate.in 2) -- (orgate.in 2);
\draw[->] (orgate.out) -- ++(1,0) node[right]{$F$};
\node at (0,-2) {A2A1A0 = 变量A,B,C};
\node at (4,-2) {OR门输入=Y1,Y3,Y5,Y7};
\end{tikzpicture}
\end{center}

\textbf{【VHDL】}
\begin{lstlisting}
-- 例化译码器
dec: decoder_3to8 port map(A=>ABC, Y=>Y);
-- OR门
F <= Y(1) or Y(3) or Y(5) or Y(7);
\end{lstlisting}

\textbf{【推广：实现任意3变量函数】}
\begin{center}
\begin{tabular}{|l|l|}
\hline
\textbf{函数} & \textbf{连接} \\
\hline
$F=m0+m2+m4+m6$ & Y0+Y2+Y4+Y6 \\
$F=m0+m1+m2+m3$ & Y0+Y1+Y2+Y3 \\
$F=A$ (即m4+m5+m6+m7) & Y4+Y5+Y6+Y7 \\
$F=\overline{A}$ & Y0+Y1+Y2+Y3 \\
\hline
\end{tabular}
\end{center}

\textbf{【如果函数用NAND实现】}
用NAND门替代OR门（德摩根定理）。此时输出是反相的，但功能等价。

\begin{examfocus}
\textbf{必考！用3-8译码器+门电路实现任意3变量逻辑函数。}
步骤：(1)将函数写成最小项之和 (2)对应Y输出用OR门连接
\end{examfocus}

\end{detailbox}


\section{变式3：用3-8译码器构建4-16译码器（片选级联）}

\begin{detailbox}[完整教学]

\textbf{【问题】}只有3-8译码器，如何得到4位地址→16个输出的4-16译码器？

\textbf{【方案】}两片3-8译码器+最高位地址A3做片选。

\textbf{【级联逻辑】}
\begin{itemize}
  \item \textbf{片0（低位）}：处理A3=0时，地址0000-0111（译码器输出Y0-Y7对应系统Y0-Y7）
  \item \textbf{片1（高位）}：处理A3=1时，地址1000-1111（译码器输出Y0-Y7对应系统Y8-Y15）
  \item A3=0→片0使能，片1禁止
  \item A3=1→片1使能，片0禁止
\end{itemize}

\textbf{【VHDL】}
\begin{lstlisting}
entity decoder_4to16 is
  port (A : in std_logic_vector(3 downto 0);
        Y : out std_logic_vector(15 downto 0));
end decoder_4to16;

architecture structural of decoder_4to16 is
  signal Y_low, Y_high : std_logic_vector(7 downto 0);
  component decoder_3to8 is
    port (A : in std_logic_vector(2 downto 0);
          E : in std_logic; Y : out std_logic_vector(7 downto 0));
  end component;
begin
  -- 片0：A3=0时使能
  dec_low: decoder_3to8 port map(
    A => A(2 downto 0), E => not A(3), Y => Y_low);

  -- 片1：A3=1时使能
  dec_high: decoder_3to8 port map(
    A => A(2 downto 0), E => A(3), Y => Y_high);

  Y <= Y_high & Y_low;  -- 拼接为16位输出
end structural;
\end{lstlisting}

\textbf{【通用级联公式】}
用$2^k$个$n$-$2^n$译码器构建$(n+k)$-$2^{n+k}$译码器。
高k位地址→片选，低n位地址→各片输入。

\textbf{【常见考题】}用两片74LS138构建4-16译码器。利用使能端做片选。

\end{detailbox}


\section{变式4：用译码器做地址译码}

\begin{detailbox}[完整教学]

\textbf{【应用场景】}
CPU有16位地址总线，但连接了多个外设（RAM、ROM、IO端口等）。
用译码器将地址空间划分为不同区域，产生各外设的片选信号。

\textbf{【例题】}用3-8译码器产生8个外设的片选信号。地址A15-A13做译码输入。

\textbf{【电路】}
\begin{center}
\begin{tabular}{|c|c|l|}
\hline
A15A14A13 & 选中输出 & 外设 \\
\hline
000 & Y0 & RAM1(0x0000-0x1FFF) \\
001 & Y1 & RAM2(0x2000-0x3FFF) \\
010 & Y2 & ROM (0x4000-0x5FFF) \\
011 & Y3 & IO1 (0x6000-0x7FFF) \\
100 & Y4 & IO2 (0x8000-0x9FFF) \\
\hline
\end{tabular}
\end{center}

\textbf{【核心思想】}高位地址→译码→片选。低位地址→片内寻址。

\end{detailbox}


\section{变式5：显示译码器（BCD→7段）}

\begin{detailbox}[完整教学]

\textbf{【问题】}3-8译码器每次只选中1个输出。但7段数码管需要同时点亮多个段（如数字"8"需要7段全亮）。

\textbf{【区别】}
\begin{itemize}
  \item 3-8译码器：$n$输入→$2^n$输出，输出是\textbf{互斥}的（恰好一个为1）
  \item BCD-7段译码器：4输入→7输出，输出\textbf{不互斥}（可以同时多个为1）
\end{itemize}

\textbf{【详见第三章BCD译码器】}BCD-7段译码器本质上是7个独立的组合逻辑函数，不是简单的"最小项展开"。

\begin{examfocus}
\textbf{关键区别}：
\begin{itemize}
  \item 3-8译码器：1-of-N选择（互斥输出）→ 用AND门
  \item BCD-7段：任意多段同时亮（非互斥）→ 用ROM/查找表
\end{itemize}
不搞混这两种"译码器"是考试的基本要求。
\end{examfocus}

\end{detailbox}


\section{译码器变式总结}

\begin{examfocus}[译码器5大变式速查]
\begin{center}
\begin{tabular}{|l|l|l|}
\hline
\textbf{变式} & \textbf{核心考点} & \textbf{常考题型} \\
\hline
74LS138风格 & 低有效+3使能端 & 写VHDL, 判断使能条件 \\
译码器+OR≈任意函数 & 最小项之和 & 给定函数→画电路 \\
级联构建更大译码器 & 高位地址做片选 & 74LS138×2→4-16 \\
地址译码 & 地址空间划分 & 存储器片选设计 \\
vs BCD-7段 & 互斥vs非互斥输出 & 概念辨析 \\
\hline
\end{tabular}
\end{center}
\end{examfocus}
